\section{Algoritmo universal de la Ley General de Masas}
\Needspace{7\baselineskip}
\subsection{Dominio del lector celular: extensión superviviente, ruta, registro y firma}
La entrada contiene únicamente:

\[
 (\mathrm{APP},\mathrm{TRIT},\mathrm{TPK},
 \mathcal R_{324},\mathcal L_{108},
 \text{representación},\text{acoplamiento},
 \text{carta dimensional común}).
\]
No contiene masa observada, residual, coeficiente elegido por aproximación ni tabla externa.

\Needspace{7\baselineskip}
\subsection{Construcción functorial de la realización material: fibra de rutas, entrelazador, proyector de coincidencia, operador espectral y escalarización}
Para cada estado físico \(X\):

\begin{enumerate}
\item construir su descriptor de carga, generación, sabor, color, espín, estadística,
orientación y composición;

\item formar la fibra de rutas compatibles;
\item calcular el espacio de entrelazadores
\end{enumerate}
   \[
   \mathscr I_X=
   \operatorname{Hom}_{G_X}(V_X,\ell^2(F_X));
   \]
\begin{enumerate}
\item si \(\dim\mathscr I_X=0\), descartar esa identificación;
\item construir el proyector mediante promedio de Reynolds y factor polar;
\item reconstruir los registros genealógicos de todas las rutas del soporte;
\item calcular firma, \(K_D\), \(K_\Omega\) y
\(\mathcal F_{\rm tower}^{2w}\);

\item aplicar el puntero de lectura del acoplamiento, el equalizador o la medida
espectral conjunta;

\item construir \(\widehat M_X\);
\item reducir a escalar sólo si la varianza es cero o existe un polo aislado;
\item para compuestos, aplicar antes \(\mathfrak C_{12}\);
\item para anchuras, calcular el semigrupo de supervivencia;
\item multiplicar por la sección dimensional común;
\item congelar la salida;
\item abrir la comparación externa.
\end{enumerate}
\Needspace{7\baselineskip}
\subsection{Existencia de la realización por fibras}
La multiplicidad de una representación \(\pi_X\) dentro de una fibra se calcula por

\[
 \dim\mathscr I_X=
 \frac1{|G_X|}
 \sum_{g\in G_X}
 \overline{\chi_{\pi_X}(g)}
 |\operatorname{Fix}_{F_X}(g)|.
\]
Si es positiva, existe un entrelazador. Si es mayor que uno, el resultado es un haz de realizaciones y el acoplamiento debe seleccionar un estado o conservar la medida mixta. Esta fórmula reemplaza una búsqueda nominal por una condición algebraica finita sobre las 324 rutas.

\Needspace{7\baselineskip}
\subsection{Igualador espectral de los lectores \(K_D\) y \(K_\Omega\) mediante \(Q_{\mathrm{eq}}=\mathbf1_{\{0\}}(\widehat L_D-\widehat L_\Omega)\)}
Sean

\[
 L_r(\Gamma)=
 \kappa_r(\Gamma)\log R_{\rm act}
 +\sum_{h\in\mathfrak S}\eta_h^r(\Gamma)s_h,
 \qquad r\in\{D,\Omega\}.
\]
El proyector de coincidencia es

\[
 Q_{\rm eq}=
 \mathbf1_{\{0\}}(\widehat L_D-\widehat L_\Omega).
\]
Una realización admite lectura común exactamente cuando

\[
 P_X\le Q_{\rm eq}.
\]
Si existe una involución del acoplamiento que intercambia ambos lectores, se usa \(L_{\rm sym}\). En otro caso, se publica su medida espectral conjunta.

\Needspace{7\baselineskip}
\subsection{Codominio del lector celular: carácter adimensional, torre y realización de masa}
La salida puede ser:

\[
 \begin{cases}
 m_X\in\mathcal L_M,&\text{masa escalar},\\
 \operatorname{Spec}(\widehat M_X),&\text{multiplete},\\
 z_X=M_X-i\Gamma_X/2,&\text{resonancia},\\
 0\in\mathcal L_M,&\text{sección nula},\\
 (\mathcal C,F,R,\Delta E),&\text{sector topológico},\\
 (L,\operatorname{Obs}_L),&\text{sección virtual}.
 \end{cases}
\]
Esta diversidad de codominios es parte de la ley; no una excepción.

\Needspace{7\baselineskip}
